\chapter{Conclusions and future work}
\label{ch:conclusions}

\begin{quote}\itshape
This chapter in brief: what was built, what the pre-registered evidence says, what an
adopting organization should do differently on Monday, and what remains honestly open.
\end{quote}

\keyidea{What transfers off this simulator is not the absolute numbers but the rankings, mechanisms and geometry --- which method wins where and why, and the observability structure any CFM56-class cockpit set shares. Every operational claim in this document is of that second, transferable kind.}


\section{What this study answers, and how}
\label{sec:questions-answered}

Before the narrative summary, the ledger: every question the study set out to answer, its
answer, and where the evidence lives. The negatives and the partial results are listed with
the same weight as the confirmations --- an honest ledger is one where refutations are as
visible as wins.

\begin{center}\footnotesize
\setlength{\tabcolsep}{4pt}
\textbf{Scientific questions (pre-registered hypotheses).}\\[2pt]
\begin{tabular}{@{}>{\raggedright\arraybackslash}p{0.30\textwidth}>{\raggedright\arraybackslash}p{0.44\textwidth}>{\raggedright\arraybackslash}p{0.18\textwidth}@{}}
  \toprule
  Question & Answer & Where \\
  \midrule
  Does AI detect faults earlier than trend monitoring? & \textbf{Yes}, once tuned: 12$\times$
    earlier at equal false alarms (recall 0.48 vs 0.13) & \cref{sec:res-h1} (H1) \\
  Does AI isolate confusable faults better? & \textbf{No}: 31\,\%=31\,\%; the wall is
    informational, not algorithmic & \cref{sec:res-h2} (H2) \\
  Does AI predict remaining life better? & \textbf{Yes vs practice} ($2.3$--$4.4\times$);
    \textbf{vs the best classical method, late life only}: $1.34\times$ at $90\,\%$
    ($p=0.006$), a tie at $70\,\%$. Robust across architectures &
    \cref{sec:res-h3h5}, \cref{sec:f8-l5,sec:f8-lrul,sec:f8-lrul-sweep} (H3) \\
  Does injecting physics into the learner help? & \textbf{Task-dependent}: no for
    prognosis, twice refuted; \textbf{yes} for mechanism attribution ($+0.139$ $R^2$,
    $p=0.0004$) & \cref{sec:res-h4}, \cref{sec:f8-l6} (H4), \cref{sec:f17-hybrid} \\
  Is the AI's uncertainty better calibrated? & \textbf{Yes, but mostly thanks to conformal,
    not to AI}: at equal coverage $1.34\times$ tighter than advanced classical, $2.3\times$
    than practice --- not the frozen $6\times$ & \cref{sec:res-h3h5} (H5),
    \cref{sec:f-uq} \\
  Can observability be bought without new hardware? & \textbf{Partly}: report-schedule design
    $+79$\,\%; learned fusion triples classical & \cref{ch:tomography} (H7) \\
  Can we certify what a diagnosis honestly knows? & \textbf{Yes}, validated against truth
    ($\rho=0.70$) & \cref{ch:breakthrough} (H10) \\
  Is the isolation wall instantaneous, or does it hold over a history? & \textbf{Instantaneous}:
    3 of 5 degradation mechanisms are attributable from trajectory shape, \code{hot\_section}
    at $R^2\,0.78$ & \cref{sec:f13-mech} \\
  Does that make learning the right tool? & \textbf{Only sometimes}: it wins on shapes that
    resist parameters, loses on those that do not (2 of 5) & \cref{sec:f13-mech} \\
  Can we tell a lying sensor from a failing engine? & \textbf{No} --- both families at chance;
    the ICM is surjective, so a real fault also fakes a bias & \cref{sec:f15-instrument} \\
  Does the certificate work as an injection channel? & \textbf{Not shown}: the control
    arm fires too, so the gain is channel count. A $C-B$ residual of $+0.053$ is
    post-hoc only & \cref{sec:f17-hybrid} \\
  Is a bidirectional network's edge the direction or the extra parameters? &
    \textbf{The parameters}: at equal budget bidirectionality buys nothing on either
    task; at doubled budget it buys $94$\,cy of RUL error and still nothing on
    attribution & \cref{sec:f18-bidir} \\
  Would transients break the confusable pair? & \textbf{Not through operating points}: the
    whole calibrated envelope gives $1.38^{\circ}$, below cruise alone. True dynamics are
    priced ($0.7$ whitened units), not ruled out & \cref{sec:gate-t} \\
  \bottomrule
\end{tabular}
\end{center}

\begin{center}\footnotesize
\setlength{\tabcolsep}{4pt}
\textbf{Operational questions (what an organization can act on).}\\[2pt]
\begin{tabular}{@{}>{\raggedright\arraybackslash}p{0.30\textwidth}>{\raggedright\arraybackslash}p{0.44\textwidth}>{\raggedright\arraybackslash}p{0.18\textwidth}@{}}
  \toprule
  Question & Answer & Where \\
  \midrule
  Where is AI \emph{not} needed? & Event-shaped faults --- the traditional stack wins at
    near-zero cost & \cref{ch:trad}, \cref{ch:cases} \\
  Which task justifies the AI investment? & Prognosis, with calibrated intervals ---
    stated against both baselines & \cref{ch:results}, \cref{sec:f8-lrul} \\
  What can today's sensors know? & Exactly the identifiable subspace the certificate maps &
    \cref{ch:breakthrough} \\
  What does textbook GPA deliver here? & A quarter of the fault recovered, $85\,\%$ smeared,
    1 parameter of 10 resolved --- before any learning & \cref{ch:gpa-practice} \\
  Is a large RUL error the method's fault or the future's? & Both, by life stage: $87\,\%$
    irreducible early; a $4\times$ reducible gap late & \cref{sec:bt-prognostic-floor} \\
  How many unscheduled removals become scheduled? & AI nets $35$--$50\,\%$ vs traditional
    $0$--$25\,\%$; capped early by the floor & \cref{sec:ops-conversion} \\
  Which sensor rescues which fault? & Station probes: HPC efficiency $45\times$; isolation
    0.15$\to$0.92 & \cref{sec:bt-honesty}, \cref{sec:f8-lh2} \\
  Would quieter sensors of the same set help? & \textbf{No}: halving every $\sigma$ leaves
    confusable isolation at $0.31$--$0.38$; buy new stations, not better probes &
    \cref{sec:noise-sweep} \\
  Can a virtual sensor or clever algorithm substitute a physical one? & \textbf{No} --- it
    adds no information (rank unchanged) & \cref{sec:f8-lh2} \\
  How much will a compressor wash recover? & Estimable from trajectory shape, $R^2=0.86$ &
    \cref{sec:f8-l4} \\
  What is it worth economically? & Median \$3\,M/yr for 100 engines, wide interval, honest
    downside & \cref{ch:economics} \\
  \bottomrule
\end{tabular}
\end{center}

\begin{center}\footnotesize
\setlength{\tabcolsep}{4pt}
\textbf{Method questions (how credible science is done here).}\\[2pt]
\begin{tabular}{@{}>{\raggedright\arraybackslash}p{0.30\textwidth}>{\raggedright\arraybackslash}p{0.44\textwidth}>{\raggedright\arraybackslash}p{0.18\textwidth}@{}}
  \toprule
  Question & Answer & Where \\
  \midrule
  How to compare AI and traditional without cheating? & Identical data, metrics, splits and
    tuning budget, frozen by pre-registration & \cref{sec:fair-design} \\
  How to trust a synthetic dataset? & Three audits and a difficulty gate that is allowed to
    fail (and did) & \cref{sec:audits} \\
  How to test your own premise? & Perturb the calibration inside its tolerance and re-read the
    verdict: robust at the core, not at the fringe & \cref{sec:f-icm} \\
  How to build the benchmark from public data alone? & Certified TCDS/EEDB measurements
    $\rightarrow$ calibrated twin & \cref{ch:model}, \cref{ch:data} \\
  How to prove an uncertainty statement is honest? & Validate it against ground truth ---
    only a synthetic-truth benchmark can & \cref{sec:bt-honesty} \\
  How to know a null result is not a broken pipeline? & Give the method a case where physics
    says the information cannot exist, and check it finds nothing & \cref{sec:f15-instrument} \\
  Does a published deep-learning PHM result survive this bar? & It replicates, but its
    architecture ranking is inside seed noise and it fails to train half the time &
    \cref{sec:f14-ext} \\
  \bottomrule
\end{tabular}
\end{center}

\section{What was built}

From nothing but certified public data, this work produced: a calibrated thermodynamic twin
of the CFM56-7B26 whose fuel-flow predictions sit within 3\,\% of certification test-bed
measurements (C1's foundation); a full classical gas-path-analysis characterization of that
engine --- signature atlas, recovery, smearing and detectability limits per sensor set
(\cref{ch:gpa-practice}); \textbf{SynCFM56}, a 100-engine run-to-failure fleet with
complete per-flight ground truth in airline snapshot format, gated by three audits that
failed and were fixed twice (C1); a traditional EHM pipeline built to industrial shape and a
matched AI suite consuming identical inputs (C2's fairness); conformal uncertainty for
prognosis (C4); the Physics-Consistency Score machinery (C5); the sensor-set observability
analysis (C6's core); an external-benchmark ranking check (C7); and a pre-registered,
Holm-corrected adjudication of five hypotheses under a symmetric 50-trial tuning budget (C2)
--- all reproducible on one desktop machine in minutes and documented to be readable without
prior specialization.

\section{What the evidence says}

Three confirmations, two refutations (\cref{tab:verdicts}), and the split is the message:

\begin{itemize}
  \item \textbf{Learning wins where patterns accumulate}: prognosis ($2.3$--$4.4\times$
    lower confirmatory RUL error against fielded practice, ranking replicated on C-MAPSS) and
    --- after honest tuning --- detection timing ($12\times$ earlier at equal false alarms).
    But the margin must be quoted at the right scale, and the right scale is unflattering:
    against the best \emph{advanced} classical prognostic the AI leads $1.34\times$ at
    $90\,\%$ of life and \emph{ties} at $70\,\%$ (\cref{sec:f8-lrul-sweep}), and its
    interval advantage falls from $6\times$ to $1.34\times$ once both sides are conformalized
    alike (\cref{sec:f-uq}) --- most of that $6\times$ was the calibration method, which is
    model-agnostic and free to the traditional side. The defensible claim is narrow:
    \emph{learning buys the late-life shop-visit horizon}. Quoting the against-practice factor
    alone would be the straw-man comparison this project exists to criticize.
  \item \textbf{Physics wins where information is scarce, and physics injection is
    task-dependent}: the ICM rule matched or beat the learner on confusable isolation, and
    bolting physics-derived features onto a learner made prognosis \emph{worse} twice (H4, L6)
    --- yet the same injection makes mechanism attribution decisively \emph{better}
    ($+0.139$ $R^2$, $p = 0.0004$, \cref{sec:f17-hybrid}). The reconciliation is that remaining
    life is already carried by the EGT-margin channel the learner sees, while ``which mechanism''
    is exactly what a health-state estimate encodes, smeared or not. ``Physics-informed'' is
    neither a free win nor a dead end; the task decides. Physics' real contribution was \emph{diagnosis of the limits}: the
    influence-matrix geometry predicted, before any learning ran, exactly where every method
    would fail.
  \item \textbf{The wall is a wall in \emph{time}, not only in geometry --- and that bounds
    it}: the confusable pair is unrecoverable at a snapshot, yet three of five degradation
    \emph{mechanisms} are attributable from the trajectory, \code{hot\_section} --- which loads
    that very pair --- at $R^2 = 0.78$ (\cref{sec:f13-mech}). The rule that emerges from four
    independent tests is \emph{trajectory shape separates what has shape}: a break-in knee and
    a wash sawtooth are readable, two plain-linear mechanisms and a slow sensor-bias ramp are
    not (\cref{sec:f15-instrument}). So the operationally useful question --- which module is
    consuming this engine --- has an answer the instantaneous inverse problem does not.
  \item \textbf{Sensors bound everyone}: the confusable wall (H2) is informational, not
    algorithmic. The extended gas-path set quintuples the worst signature angle; no model
    recovers what probes never measured --- and halving the noise of the \emph{existing}
    probes does not help either (\cref{sec:noise-sweep}), which makes the recommendation
    specific: new stations, not better probes at the old ones.
  \item \textbf{Learning degrades gracefully, thresholds fall off a cliff}: across a factor of
    eight in sensor noise the AI's remaining-life error is flat and its detection recall falls
    by a third, while the tuned classical detector goes from $0.48$ recall to zero and the
    margin extrapolation loses $40\,\%$ of its accuracy. Robustness to instrumentation quality,
    not peak accuracy, may be learning's most transferable advantage here.
\end{itemize}

\paragraph{The refutations did more work than the confirmations.} Worth stating plainly,
because it is the clearest argument for running a study this way. The three confirmations
said roughly what the field expected: learning predicts remaining life better, detects
earlier once tuned, quantifies uncertainty better. Useful, but not surprising. The two
refutations are what changed the recommendations. H2 converted an algorithm question into a
purchase order --- \emph{buy probes at new stations} --- which is the single most actionable
sentence in this document. H4 removed a popular technique from the shortlist, twice, saving
an adopter the effort of discovering it themselves. A study whose hypotheses all confirm has
usually either asked easy questions or moved its own goalposts.

\section{Guidance for adoption}

The five-point checklist of \cref{sec:res-industry} stands as the operational summary:
prognosis first, tuned change-detection second, sensors before models for isolation,
distrust of untuned comparisons, and pre-registration as a procurement demand. To it the
case studies add a sixth: \textbf{protect the sensor estate} --- the drifting thermocouple
of Case C corrupts every consumer downstream, learned or classical.

\section{What it is worth}

\Cref{ch:economics} puts a deliberately cautious number on it: for a 100-engine 737NG
operator, the one mechanism the results support --- less-biased prognosis converting
unscheduled removals to scheduled --- plausibly saves low single-digit millions per year
(median \$3\,M, p10 \$0.9\,M), with an honest structural chance of nothing. The larger,
unquantified prize is the diagnostic instrumentation the certificate would justify buying.

\section{Honest limitations}
\label{sec:limitations-conclusions}

The world here is synthetic and linearized (audited to the noise floor, but still a model
of a model); the mission mix is two snapshot conditions with scatter; the EGT is a
station-4.5 proxy; chronic labels correlate with age by construction; the AI family is one
compact architecture per task, not the state of the art; and H4's verdict covers one of
three planned hybrid mechanisms.

One limitation was found while trying to build on the fleet and deserves naming, because it
bounds what \cref{sec:f13-mech} can be used for. \textbf{The fleet is homogeneous in
workscope.} Every engine draws all five degradation mechanisms with lognormal severity
multipliers, and \code{hot\_section} carries the largest rate, so at $95\,\%$ of life it is the
dominant contributor to lost EGT margin in \textbf{100 of 100 engines} (mean share $0.53$,
against $0.16$ for \code{clearance} and $0.06$ for \code{fouling}). A shop-visit question of
the form ``which module will this engine need'' therefore has a constant answer here, and any
hit-rate metric built on it is scored against a $100\,\%$ majority class. Real fleets are not
like this --- module-specific failures, foreign-object damage and build-standard variation
produce workscope heterogeneity this generator does not model. The consequence is specific:
\cref{sec:f13-mech}'s result is a statement about \emph{identifiability} --- how much of each
mechanism is recoverable from a trajectory --- and must not be read as a claim that maintenance
planning can be predicted from it on this benchmark. The operationally meaningful residue is
\emph{how much} rather than \emph{which}, and that question already has its answer in
\cref{sec:f8-l4}. None of these caveats is hidden in fine print --- each is
in the decision register or the relevant chapter.

\section{The capstone: earned-confidence diagnosis}

The project's furthest reach (\cref{ch:breakthrough}) turns its two rarest assets --- a
calibrated differentiable twin and complete ground truth --- into something the field has
wanted and never had: a per-engine certificate of \emph{what gas-path diagnosis can and
cannot know}, proven honest against the true component state ($\rho = 0.70$ between certified
precision and actual error, $p = 0.025$). It converts smearing from a hidden failure into an
explicit, trustworthy, per-direction confidence statement, and the same Fisher-information
object tells an operator which sensor rescues which fault ($45\times$ for HPC efficiency
from station probes) --- unifying the H2 wall and the F7 schedule design under one rigorous
instrument. The honest boundary is stated too: the CRB earns the region's validated anisotropic shape,
and a disclosed split-conformal factor earns its absolute scale (restoring a 90\,\% coverage
guarantee). This is the contribution most likely to matter outside
this document, because it is the one no real-data study can validate.

A second certificate (\cref{sec:bt-prognostic-floor}) does the analogous job for prognosis.
Conditioning on each engine's \emph{true} current health state removes all epistemic
ignorance, exposing an aleatoric floor on remaining-life error that no method can beat: at
mid-life $87\,\%$ of the RUL spread is irreducible, yet a $4\times$ reducible gap opens late
in life. It answers the question every operator should ask before funding a better prognostic
--- is the error the method's or the future's --- with a ground-truth-validated number and the
life stage at which investment actually pays.

\section{The F7 extension}

After the main study closed, the operating-point tomography phase (\cref{ch:tomography})
converted the isolation stalemate into two new results under a second pre-registered
freeze: the report schedule as an observability design variable ($+79\,\%$ separability
from one procedural report condition), and learned fusion of physics projections tripling
classical multi-point stacking (0.65 vs 0.17) with calibrated, physics-tracking ambiguity
sets --- plus one honestly refuted threshold hypothesis. The limitations-overcoming program turns every declared limitation into a research line
(\cref{ch:overcoming}). Three are complete: L1 replaced the linearized generator's physics
with a differentiable neural twin 13--19$\times$ more faithful (\cref{fig:surrogate}); L2
regenerated the fleet through it (SynCFM56 v2, 3--7$\times$ closer to full physics,
difficulty gate re-passed); and L6 used v2 to re-open H4 with a twin-residual feature and
\textbf{refuted it a second, independent way} (\cref{fig:h4v2}). Four further lines followed,
each pre-registered: L5 showed the RUL advantage is \emph{architecture-robust} (a TCN and a
Transformer beat the traditional baseline as well as the GRU, \cref{sec:f8-l5}); L4
established that the \emph{recoverable fraction} of lost margin is estimable from trajectory
shape ($R^2 = 0.86$, a wash-planning tool, \cref{sec:f8-l4}); L7 built the augmented-state
Kalman for sensor drift, which tracks the drift but cannot un-corrupt the cockpit diagnosis
(\cref{sec:f8-l7}); and L9 could only partially validate the PCS metric, because no available
classifier is strongly enough physical to test it against (\cref{sec:f8-l9}). The pattern is
the project's in miniature: two clean positives that extend capability, two partial results
that reinforce the observability wall. The three heaviest remaining lines (real-station EGT,
N-CMAPSS DS02, map-shape calibration) are characterized with hypotheses and gates in the
plan.

\section{Boundaries of use}
\label{sec:boundaries}

This work produces diagnostic and prognostic estimates, and one of its lines recommends
\emph{which data to collect next} (\cref{ch:tomography,sec:bt-honesty}). Both sit near
regulated activity, so the boundary is stated explicitly rather than left to be inferred; the
full statement is \code{docs/safety-case-boundaries.md}.

\textbf{What this is:} an offline engineering study, a benchmark and measurement method, and a
decision-support instrument for engineering judgement --- which health directions a sensor set
can resolve, how much of a remaining-life error is irreducible, which added measurement resolves
the most.

\textbf{What it is not, and must not become without a case of its own:} it does not control the
engine, does not interface with the FADEC or any certified function, does not design new
operational manoeuvres (a recommended \emph{observation} selects among conditions that are
already authorised and already flown), does not modify limits or protections, and does not issue
airworthiness determinations.

\textbf{The rule that matters most:} \emph{a statistical prediction is never presented as a
certified conclusion.} Three of this document's own findings exist to bound what may be claimed
--- $87\,\%$ of mid-life remaining-life spread is irreducible, so a point estimate is not a
failure date (\cref{sec:bt-prognostic-floor}); the confusable pair is unresolvable from cockpit
sensors, so naming one component is not evidence the other is healthy (\cref{sec:res-h2});
and a drifting sensor cannot be told from real degradation, so every diagnosis here is
conditional on instrument health the method cannot verify (\cref{sec:f15-instrument}).

Every output is advisory to a qualified engineer, traceable to the script, pre-registration tag,
data version and seeds that produced it. Before any operational use, four prerequisites are
listed in the boundaries document, and none has been done --- validation on real engine data,
a verification case covering the silent-collapse and seed-instability failure modes this project
itself documented (\cref{sec:f14-ext,sec:f18-bidir}), an assessment against the applicable
regulatory framework, and human-factors work on presenting an interval so it is not read as a
certainty.

\section{Future work}

\subsection{The one line that must be run}
\label{sec:future-f12}

The re-audits of \cref{sec:f8-lrul-sweep,sec:f-uq} left the AI case honest and small: a
$1.34\times$ late-life prognostic edge over a competent classical method, a tie at mid-life, a
detection win traceable to window search, a refuted isolation claim and a twice-refuted
hybrid. That is a defensible result. It is not a mechanism, and this document should not
pretend otherwise.

The open question worth the project's remaining effort is therefore not \emph{how much better
is AI}, which is now measured, but \emph{is there anything a learner can do here that a
competent classical method structurally cannot}. There is a candidate, and it follows from
this document's own findings. Classical gas-path analysis closes the seven unconstrained
dimensions of an underdetermined inverse problem with a fixed regularizer --- a Gaussian,
near-diagonal, engine-independent, life-invariant prior. The real fleet prior is correlated
across components, non-Gaussian, non-stationary and engine-specific. \textbf{The one thing a
learner can structurally do is carry that prior into the measurement's null space.} Under this
reading, H2 failed because the learner was handed the same per-snapshot framing as GPA, in
which there is no prior to exploit, and prognosis wins late because trajectory shape --- a
fleet-level statistic --- only becomes informative once degradation dominates. Both are
predictions of the claim, not post-hoc stories.

What makes it testable \emph{here}, and effectively nowhere else, is that the certificate of
\cref{ch:breakthrough} supplies an \emph{absolute} reference rather than a baseline: the
Cram\'er--Rao bound is the precision floor for any unbiased estimator using one engine's data
alone. It cannot be under-tuned or strawmanned --- the objection that shrank H3 and H5. A
biased or Bayesian estimator is entitled to beat it; what is normally impossible is
\emph{verifying} that it did, because verification needs the true component state. So an
estimator that beats the CRB in a certified-unidentifiable direction is provably importing
information from outside that engine's measurement, and the size of the violation is the size
of the import --- per engine, per direction, in physical units. Summed as
$\frac{1}{2}\log_2(\det\Sigma_{\text{CRB}} / \det\Sigma_{\text{achieved}})$ it is a single
number in bits: \emph{how much of this diagnosis came from the sensors, and how much from the
fleet}.

The proposal (\code{docs/f12-proposal.md}) freezes five gates, each allowed to fail, and the
two that matter most are the ones that can cost the project its claim. A classical steelman
receives the \emph{same} empirical fleet prior as its a-priori covariance; if it recovers most
of the gain, the honest headline becomes ``the advantage is the prior, and it costs twenty
lines of linear algebra, not a neural network'' --- the most useful anti-hype result available
here. And a transfer gate trains the prior on one fleet and tests it on another whose
degradation-mode mixture is deliberately mismatched, which turns this document's deepest
unexamined assumption --- that a certificate validated against its own generator's truth means
something outside that generator --- into an explicit experimental variable. A fifth gate is a
built-in fraud detector: the gain must appear \emph{where the certificate says it should},
large in unobservable directions and negligible in identifiable ones. Leakage and overfitting
produce gains everywhere; only a genuine prior produces gains shaped like the null space, and
that shape is predicted in advance by an object computed without ever looking at the
estimator. Every branch, including outright failure at the first gate, terminates in a result
worth stating.

\subsection{The remaining programme}

(1) The physics-constrained-loss mechanism (M3), now deferred behind two H4 refutations and
requiring its own pre-registration against that record; (2) the PCS evaluated on the
competent F7 isolation model (C5 validation); (3) N-CMAPSS DS02 as the heavyweight
sim-to-real extension; (4) mission-mix and transient realism in the generator, feeding the
report-schedule design of \cref{ch:tomography} with data rather than geometry; (5) the
extended-sensor fleet, turning the C6 observability analysis into a measured acquisition
trade study; (6) public release of SynCFM56 (v1.1 frozen + v2 nonlinear) with a DOI.
